% Part: incompleteness
% Chapter: representability-in-q
% Section: c-representable

\documentclass[../../../include/open-logic-section]{subfiles}

\begin{document}

\olfileid{inc}{req}{cre}
\olsection{C 中函数在 $\Th{Q}$ 中可表示}

我们需要证明 $C$ 中的每个函数都在 $\Th{Q}$ 中可表示。归根结底，我们需要展示如何为 $C$ 中的每个 $k$ 元函数 $f(x_0,\dots,x_{k-1})$ 指派一个表示它的公式 $!A_f(x_0,\dots,x_{k-1},y)$。

\begin{lem}
给定自然数 $n$ 和 $m$，若 $n \neq m$，则 $Q \Proves \num n \neq \num m$。
\end{lem}

\begin{proof}
对 $n$ 进行归纳以证明：对每个 $m$，若 $n \neq m$，则 $Q \Proves \eq/[\num n][\num m]$。

基础情形下，$n = 0$。若 $m$ 不等于 $0$，则存在某个自然数 $k$ 使得 $m = k + 1$。有公理表明 $\lforall[x][\eq/[0][x']]$。由量词公理，将 $x$ 替换为 $\num k$，可得 $\eq/[0][\num k']$。而 $\num k'$ 即 $\num m$。

归纳步骤中，假设命题对 $n$ 成立，考虑 $n+1$。令 $m$ 是任意自然数。有两种可能：$m = 0$，或存在 $k$ 使得 $m = k+1$。第一种情形如上处理。第二种情形下，假设 $n+1 \neq k+1$。那么 $n \neq k$。根据对 $n$ 的归纳假设，有 $\Th{Q} \Proves \eq/[\num n][\num k]$。有公理表明 $\lforall[x][\lforall[y][\eq[x'][y'] \lif \eq[x][y]]]$。由量词公理，得 $\eq[\num n'][\num k'] \lif \eq[\num n][\num
  k]$。由命题逻辑，可在 $\Th{Q}$ 中推得 $\eq/[\num n][\num k] \lif \eq/[\num n'][\num k']$。通过肯定前件规则，可得 $\eq/[\num n'][\num k']$，这即为所求，因为 $\num k'$ 即 $\num m$。
\end{proof}

注意，该引理本身并非什么很强的结论：本质上它只是表明 $\Th{Q}$ 能证明不同的数字表示不同的对象。例如，$\Th{Q}$ 证明了 $0'' \neq 0'''$。但要证明它普遍成立，仍需谨慎处理。还需注意，尽管我们使用了归纳法，但这是在 $\Th{Q}$ \emph{之外} 进行的归纳。

我们将能够表示零、后继、加法、乘法、等词的特征函数以及投影函数。每种情形下，相应的表示公式都十分直接；例如，零由公式
\[
y = 0,
\]
表示，后继由公式
\[
x_0' = y,
\]
表示，加法由公式
\[
x_0 + x_1 = y.
\]
表示。主要工作在于证明 $\Th{Q}$ 能够证明相关的语句；例如，断言加法可由上述公式表示，就意味着需要证明对每一对自然数 $m$ 和 $n$，$\Th{Q}$ 都证明
\[
\eq[\num n + \num m][\num {n+m}]
\]
和
\[
\lforall[y][(\eq[(\num n + \num m)][y] \lif \eq[y][\num {n+m}])].
\]

对于复合运算呢？假设 $h$ 由
\[
h(x_0,\dots,x_{l-1}) = f(g_0(x_0,\dots,x_{l-1}), \dots,
g_{k-1}(x_0,\dots,x_{l-1})).
\]
定义，其中我们已经分别找到了表示函数 $f,g_0,\dots,g_{k-1}$ 的公式 $!A_f, !A_{g_0}, \dots,
!A_{g_{k-1}}$。那么我们可以定义一个表示 $h$ 的公式 $!A_h$，只需将 $!A_h(x_0,\dots,x_{l-1},y)$ 定义为
\begin{multline*}
  \lexists[z_0\dots][\lexists[z_{k-1}][(!A_{g_0}(x_0,\dots,x_{l-1},z_0) \land
  \dots \land !A_{g_{k-1}}(x_0,\dots,x_{l-1},z_{k-1}) \land]]\\
  !A_f(z_0,\dots,z_{k-1},y)).
\end{multline*}

最后，考虑无界搜索。假设 $g(x,\vec z)$ 是正则的并且在 $\Th{Q}$ 中可表示，设其由公式 $!A_g(x,\vec
z,y)$ 表示。令 $f$ 由 $f(\vec z) = \mu x \; g(x,\vec z)$ 定义。我们想要找到一个表示 $f$ 的公式 $!A_f(\vec z,y)$。一个自然的选择是：
\[
!A_f(\vec z,y) \ident !A_g(y,\vec z,0) \land \lforall[w][(w < y
\lif \lnot !A_g(w,\vec z,0))].
\]

借助一些额外的引理（如下面的引理）可以证明，这种构造是有效的：

\begin{lem}
  对每个变元 $x$ 和每个自然数 $n$，$\Th{Q}$ 可证 $\eq[(x'
  + \num n)][(x + \num n)']$。
\end{lem}

再次值得提醒的是，这弱于断言 $\Th{Q}$ 证明 $\lforall[x][\lforall[y][(x' + y) = (x +y)']]$（而后者为假）。

\begin{proof}
照例通过对 $n$ 的归纳进行证明。基础情形下，$n = 0$，需证 $\Th{Q}$ 可证 $\eq[(x'+0)][(x+0)']$。但我们有：
\begin{eqnarray*}
& & \eq[(x' + 0)][x'] \quad \text{由公理 4} \\
& & \eq[(x + 0)][x] \quad \text{由公理 4} \\
& & \eq[(x+0)'][x'] \quad \text{由第 2 行} \\
& & \eq[(x' + 0)][(x + 0)'] \quad \text{第 1 行与第 3 行}
\end{eqnarray*}
归纳步骤中，可假设我们已在 $\Th{Q}$ 中推导出 $\eq[(x' +
  \num n)][(x + \num n)']$。由于 $\num{n+1}$ 就是 $\num n'$，需证 $\Th{Q}$ 可证 $\eq[(x' + \num n')][(x + \num
n')']$。以下等式链可在 $\Th{Q}$ 中导出：
\begin{eqnarray*}
(x' + \num n') & \eq & (x' + \num n)' \quad \text{公理 5} \\
& \eq & (x + \num n')' \quad \text{由归纳假设}
\end{eqnarray*}
\end{proof}

\begin{lem}
\begin{enumerate}
\item $\Th{Q}$ 证明 $\lnot (x < \num 0)$。
\item 对每个自然数 $n$，$\Th{Q}$ 证明
\[
x < \num {n+1} \lif (\eq[x][\num 0] \lor \dots \lor \eq[x][\num n]).
\]
\end{enumerate}
\end{lem}

\begin{proof}
我们非形式地给出结论 1 和部分结论 2 的证明（即仅提示如何构造形式推导）。

对于结论 1，根据 $<$ 的定义，我们需要在 $\Th{Q}$ 中证明 $\lnot
\lexists[y][\eq[(y'+x)][0]]$，这等价于（使用一阶逻辑的公理和规则）$\lforall[y][\eq/[(y' +
    x)][0]]$。思路如下：假设 $\eq[(y'+x)][0]$。若 $x$ 为 $0$，则有 $\eq[(y' + 0)][0]$。但由 $\Th{Q}$ 的公理 4，有 $\eq[(y' + 0)][y']$，再由公理 2 可得 $\eq/[y'][0]$，矛盾。故 $\lforall[y][\eq/[(y' +x)][0]]$。若 $x$ 不为 $0$，由公理 3，存在 $z$ 使得 $\eq[x][z']$。则有 $\eq[(y' + z')][0]$。由公理 5，得 $\eq[(y' + z)'][0]$，再次与公理 2 矛盾。

对于结论 2，对 $n$ 进行归纳。考虑基础情形 $n = 0$。此时需要证明 $x < \num 1 \lif \eq[x][\num
  0]$。假设 $x < \num 1$。由 $<$ 的定义公理，有 $\lexists[y][\eq[(y'+x)][0']]$。设 $y$ 具有该性质，于是有 $\eq[y'+x][0']$。

我们需要证明 $\eq[x][0]$。由公理 3，若 $x$ 不为 0，则存在 $z$ 使得 $x$ 等于 $z'$。则有 $\eq[(y' + z')][0']$。由 $\Th{Q}$ 的公理 5，得 $\eq[(y'+z)'][0']$。再由公理 1，得 $\eq[(y' +
    z)][0]$。但由定义，这意味着 $z < 0$，与结论 1 矛盾。
\end{proof}

\begin{explain}
我们已证明，可计算函数集可以刻画为在 $\Th{Q}$ 中可表示的函数集。事实上，该证明更具一般性。从可表示性的定义不难看出，任何扩张 $\Th{Q}$（或能解释 $\Th{Q}$）的理论都能表示可计算函数；但是反过来，在任何证明概念是可计算的推导系统中，每个可表示函数都是可计算的。因此，举例来说，可计算函数集也可以刻画为在 皮亚诺算术中，甚至在 策梅洛-弗兰克尔集合论中可表示的函数集。正如 哥德尔所指出的，这多少有些令人惊讶。我们将会看到，在可证性方面，问题对你所考虑的理论非常敏感；粗略地说，公理越强，能证明的命题就越多。然而，在极其广泛的公理理论中，可表示的函数恰恰就是可计算函数。
\end{explain}

\end{document}